\documentclass{article}
\usepackage{kotex}
\usepackage{setspace}  % 줄간격
\setstretch{1.15}
\usepackage{enumitem}  % 목록 들여쓰기 
\setlist[itemize]{leftmargin=*, align=left}
\setlist[description]{
  font=\normalfont,
  labelwidth=4cm,
  leftmargin=4cm,
  labelsep=0cm,
  align=left
}
\usepackage{titlesec}  % 섹션 스타일 조정
\titleformat{\section}
  {\normalfont\Large\bfseries}
  {}
  {0pt}
  {}
  [\vspace{0.3em}\titlerule]

\titlespacing*{\section}
  {0pt}
  {1.5em}
  {0.8em}
\newcommand{\descriptionrule}{%
  \item[]\vspace{-0.6em}
  \noindent\hspace*{-\leftmargin}%
  \rule{0.6\linewidth}{0.4pt}
  \vspace{-0.4em}
}              
\usepackage[a4paper, margin=2.5cm]{geometry}
\usepackage{fancyhdr}                 
\pagestyle{fancy}
\fancyhf{}                             
\fancyfoot[L]{\nouppercase{언어습득 (2026학년도 2학기)}} 
\fancyfoot[R]{\thepage}                
\renewcommand{\headrulewidth}{0pt}    
\renewcommand{\footrulewidth}{0.4pt}   

\begin{document}

\title{언어습득 (2026학년도 2학기)}
\author{\textit{written by using} \LaTeX}
\date{\today}
\maketitle
\thispagestyle{fancy}
\setcounter{secnumdepth}{0}
\bigskip

\section{0. 강좌정보}
\medskip
교과목번호: M1246.001200(001)\\
개설학과: 서울대학교 인문대학 언어학과\\
담당교수: 황효성\\
수업시간: 월/수 12:30--13:45\\
수업장소: 83-604

\section{1. 수업목표}
\medskip
이 강좌에서는 출생부터 취학 전 아동의 언어 습득을 개괄적으로 살펴본다. 인간의 언어 습득 능력을 설명하는 이론을 소개하고, 영유아와 아동의 언어 습득 과정을 말소리, 어휘, 의미, 문법, 화용 단계로 나누어 습득의 자세한 과정과 원리에 대해 탐구한다. 일반적인 발달에 초점을 맞추되, 이중 언어 습득과 언어 장애에 대해서도 살펴본다. 실제 언어 습득 데이터를 분석하는 방법을 배우고, 언어 습득 분야의 연구 방법에 대한 이해를 함양한다. 이 과정이 끝나면 수강생들은 아동 언어 습득의 내용, 단계 및 과정, 이를 연구하는데 사용되는 패러다임과 이를 설명하려는 주요 이론에 대해 폭넓게 이해하게 된다.

\section{2. 교재 및 참고문헌}
\medskip
\subsection{교재}
『Language development』(Erika Hoff, 2014.)

\newpage
\section{3. 주차별 강의계획\protect\footnotemark}
\footnotetext{세부 일정이 공개되지 않은 관계로 임의로 기입하였음. 모든 일정은 추후 변동 가능}
\medskip
\begin{description}
    \item[1주차(9/2)] 강의 소개
    \item[2주차(9/7, 9/9)] 언어 습득 이론
    \item[3주차(9/14, 9/16)] 언어 습득의 생물학적 기초
    \item[4주차(9/21, 9/23)] 언어 발달의 기초
    \item[5주차(9/28, 9/30)] 언어 습득 연구 방법
    \item[6주차(10/7, 10/12)] 영아의 말소리 지각 및 산출
    \item[7주차(10/14, 10/19)] 음운 발달 및 단어 분절
    \item[8주차(10/26, 10/28)] 초기 어휘 발달 및 어휘와 의미
    \item[9주차(11/2, 11/4)] 문법 발달 (1)
    \item[10주차(11/9, 11/11)] 문법 발달 (2)
    \item[11주차(11/16, 11/18)] 화용 발달
    \item[12주차(11/23, 11/25)] 이중언어 발달, 수어 발달
    \item[13주차(11/30, 12/2)] 아동 언어 장애
    \item[14주차(12/7, 12/9)] 연구계획서 발표(영상) 및 종합 복습
    \item[15주차(12/14)] 기말고사
\end{description}

\section{4. 평가방법}
\medskip
\begin{description}
    \item[\textbf{성적부여방식}] 상대평가(A--F)\footnote{S/U 선택가능}
    \item[\textbf{출석(5\%)}] 수업일수의 1/3을 초과하여 결석하면 성적은 ``F'' 또는 ``U''가 됨\\(담당교수가 불가피한 결석으로 인정하는 경우는 예외로 할 수 있음)
    \item[\textbf{과제(50\%)}] 
    \item[\textbf{중간(0\%)}] 미시행
    \item[\textbf{기말(20\%)}] 
    \item[\textbf{수시평가(20\%)}] 
    \item[\textbf{태도(5\%)}]  
    \descriptionrule
    \item[\textbf{합계}] 100\%
\end{description}

\newpage
\section{5. 생성형 AI 도구 활용 방침}
\medskip
\begin{itemize}
    \item 본 강의에서는 생성형 AI를 학습 보조 도구로 활용할 수 있습니다. 영어로 된 연구 논문을 읽고 이해하는 데 번역·요약 보조 도구의 사용을 허용하며, 과제의 아이디어 및 주제 탐색, 자료 검색, 표현 다듬기 등의 목적으로도 사용할 수 있습니다.
    \item 다만 과제의 핵심 내용은 수강생 본인의 것이어야 합니다. 비판적 요약문의 분석과 비판, 연구계획서의 연구 질문과 설계는 스스로 논문을 읽고 사고한 결과여야 하며, 논문을 직접 읽지 않고 요약을 생성하거나 계획서 초안 자체를 생성하는 것은 허용되지 않습니다.
    \item 생성형 AI를 과제 수행에 사용한 경우, 사용 도구명과 사용 과정(무엇을 위해 어떻게 사용했는지)을 과제 말미에 간략히 명시해야 합니다.
\end{itemize}

\section{6. 정원 외 신청 수용 여부}
\medskip
최대 0명

\section{7. 수강생 참고사항}
\begin{itemize}
    \item 교재를 꼭 구매할 필요는 없으며, 강의 노트와 읽기 자료는 수업 전 eTL에 게시함.
    \item 매주 관련 연구 논문을 함께 읽고 논의하며, 수강생은 수업 전날까지 해당 논문에 대한 질문을 하나씩 eTL에 게시함.
    \item 배정에 따라 논문의 일부를 맡아 발표하며, 일부 차시에는 조별 논문 분석 활동과 아동 언어 데이터 실습을 함께 진행함.
    \item 비판적 요약문을 총 2회 개별 제출하고, 학기 중 퀴즈를 3회 실시함.
    \item 기말 과제로 연구계획서를 작성하며, 발표는 영상 제출과 동료 피드백으로 진행함.
    \item 구체적인 일정과 진행 방식은 강의 첫날에 공지함.
\end{itemize}

\newpage
\section{8. 장애학생 지원사항}
\medskip
\subsection{강의 수강 관련}
\medskip
\begin{itemize}
    \item 시각장애: 교재 제작(디지털교재, 점자교재, 확대교재 등), 대필도우미 허용
    \item 지체장애: 교재 제작(디지털교재), 대필도우미 및 수업보조 도우미 허용
    \item 청각장애: 대필 및 문자통역 도우미 활동 허용, 강의 녹취 허용
    \item 건강장애: 질병 등으로 인한 결석에 대한 출석 인정, 대필도우미 허용
    \item 학습장애: 대필도우미 허용
    \item 지적장애/자폐성장애: 대필도우미 및 수업 멘토 허용
\end{itemize}
\subsection{과제 및 평가 관련}
\medskip
\begin{itemize}
    \item 시각장애/지체장애/청각장애/건강장애/학습장애: 과제 제출기한 연장, 과제 제출 및 응답 방식의 조정, 평가 시간 연장, 평가 문항 제시 및 응답 방식의 조정, 별도 고사실 제공
    \item 지적장애/자폐성장애: 개별화 과제 제출 및 대체 평가 실시
\end{itemize}
\subsection{비고}
\medskip
본 강의를 수강하는 장애학생들에게는 이상의 지원 서비스 이외에도 장애학생 개개인의 특성과 요구에 따라, 지도교수 및 장애학생지원센터와의 상담을 통하여 적절한 수준의 지원 서비스를 제공합니다. 장애학생에 대한 지원서비스와 관련하여 문의사항이 있는 학생들은 담당교수 최계영(02-880-6163) 혹은 장애학생지원센터(02-880-8787)로 문의바랍니다.

\end{document}